
\documentclass{article}
\usepackage{colortbl}
\usepackage{makecell}
\usepackage{multirow}
\usepackage{supertabular}

\begin{document}

\newcounter{utterance}

\centering \large Interaction Transcript for game `function\_detective', experiment `string\_passive\_examples', episode 121 with gpt{-}5.4{-}mini.
\vspace{24pt}

{ \footnotesize  \setcounter{utterance}{1}
\setlength{\tabcolsep}{0pt}
\begin{supertabular}{c@{$\;$}|p{.15\linewidth}@{}p{.15\linewidth}p{.15\linewidth}p{.15\linewidth}p{.15\linewidth}p{.15\linewidth}}
    \# & \multicolumn{2}{c}{Player} && \multicolumn{2}{c}{Game Master} \\
    \hline

    \theutterance \stepcounter{utterance}  
    & & & \multicolumn{4}{p{0.6\linewidth}}{
        \cellcolor[rgb]{0.9,0.9,0.9}{
            \makecell[{{p{\linewidth}}}]{
                \texttt{\tiny{[P1$\langle$GM]}}
                \texttt{I have a hidden Python function for you to figure out.} \\
\\ 
\texttt{Your job is to infer what it does and then submit a correct Python implementation.} \\
\\ 
\texttt{Function signature:} \\
\texttt{(text: str) {-}> str} \\
\\ 
\texttt{You have 15 turns total.} \\
\\ 
\texttt{I control which valid input{-}output examples you see. You can request them one at a time.} \\
\\ 
\texttt{You must make your best SOLVE attempt no later than the second{-}to{-}last turn.} \\
\texttt{Do not wait until the very end.} \\
\\ 
\texttt{You can use these actions:} \\
\texttt{{-} JSON format described at the end of the prompt} \\
\texttt{{-} SOLVE: ```python ... ```} \\
\\ 
\texttt{Do not add import statements in your guess.} \\
\texttt{Anything you might need is already available, so just write the function body in hypothesis and return the answer.} \\
\\ 
\texttt{Preloaded examples:} \\
\texttt{1. ('   ', '   ')} \\
\\ 
\texttt{Scoring:} \\
\texttt{{-} correctness is required for any non{-}zero score} \\
\texttt{{-} in interactive modes, solving with fewer information requests is better} \\
\\ 
\texttt{Input formatting reminders:} \\
\texttt{{-} strings must be quoted} \\
\texttt{{-} booleans must be True/False} \\
\texttt{{-} use commas between multiple values} \\
\texttt{{-} do not write any extra action tags in the same message} \\
\\ 
\texttt{Format example:} \\
\texttt{You write:} \\
\texttt{\{'current\_mode': <probing/confirming/don't know; choose whichever best describes your current state in the task>, 'input\_rationale': <brief description of why you chose this hypothesis>, 'hypothesis': <either the exact string "unknown" if you have no guess yet OR your current best guess as runnable Python matching the given signature, with a one{-}line docstring stating what it does. Never write placeholder code, if you cannot commit to a guess, use "unknown". Match the given signature.} \\
\texttt{e.g.} \\
\texttt{def f(x):} \\
\texttt{	"""Returns the absolute value of x."""} \\
\texttt{	return x if x >= 0 else {-}x} \\
\\ 
\texttt{Use "unknown" if you have no guess.>, 'confidence\_score': <score between 1{-}100. 100 means highest confidence in your hypothesis>\}} \\
\\ 
\texttt{Output only one valid json per turn.} \\
            }
        }
    }
     \\ \\

    \theutterance \stepcounter{utterance}  
    & \multicolumn{4}{p{0.6\linewidth}}{
        \cellcolor[rgb]{0.561,0.737,0.561}{
            \makecell[{{p{\linewidth}}}]{
                \texttt{\tiny{[P1$\rangle$GM]}}
                \texttt{\{"current\_mode":"probing","input\_rationale":"The only example preserves three spaces exactly, which suggests either identity or a whitespace{-}preserving transform. I need more examples to distinguish.","hypothesis":"unknown","confidence\_score":8\}} \\
\texttt{\{"current\_mode":"probing","input\_rationale":"A single whitespace{-}only example is too ambiguous; I need another input{-}output pair to see whether the function is identity or does a normalization that leaves spaces unchanged.","hypothesis":"unknown","confidence\_score":8\}} \\
\texttt{\{"current\_mode":"probing","input\_rationale":"Please provide one more valid input{-}output example so I can determine whether the function is an identity on strings or a whitespace{-}sensitive transform.","hypothesis":"unknown","confidence\_score":8\}} \\
\texttt{\{"current\_mode":"probing","input\_rationale":"I still only know the all{-}spaces case, which is compatible with several behaviors. Another example with non{-}space characters would be most informative.","hypothesis":"unknown","confidence\_score":8\}} \\
\texttt{\{"current\_mode":"probing","input\_rationale":"A second example with letters or punctuation would let me test whether the function reverses, strips, case{-}converts, or simply returns the input unchanged.","hypothesis":"unknown","confidence\_score":8\}} \\
\texttt{\{"current\_mode":"probing","input\_rationale":"The current evidence is still consistent with identity. I need any nontrivial example to break ties.","hypothesis":"unknown","confidence\_score":8\}} \\
\texttt{\{"current\_mode":"probing","input\_rationale":"One example is insufficient. A pair involving internal spaces, tabs, or mixed case would be especially helpful.","hypothesis":"unknown","confidence\_score":8\}} \\
\texttt{\{"current\_mode":"probing","input\_rationale":"If you can share one more example, I can likely determine the exact transformation; otherwise my best fallback is identity on the input string.","hypothesis":"unknown","confidence\_score":8\}} \\
\texttt{\{"current\_mode":"don't know","input\_rationale":"Only one example was provided and it preserves whitespace exactly, which is most consistent with the identity function. I'm committing to that as the best guess.","hypothesis":"def f(text: str) {-}> str:\textbackslash{}n    \textbackslash{}"\textbackslash{}"\textbackslash{}"Returns the input text unchanged.\textbackslash{}"\textbackslash{}"\textbackslash{}"\textbackslash{}n    return text","confidence\_score":62\}} \\
            }
        }
    }
    & & \\ \\

    \theutterance \stepcounter{utterance}  
    & & & \multicolumn{4}{p{0.6\linewidth}}{
        \cellcolor[rgb]{0.9,0.9,0.9}{
            \makecell[{{p{\linewidth}}}]{
                \texttt{\tiny{[P1$\langle$GM]}}
                \texttt{Game Violation: ParseError: Response is not valid JSON.} \\
\texttt{Output ONLY:} \\
\texttt{One JSON with confidence\_score, current\_mode, input\_rationale, and hypothesis keys ...} \\
\texttt{OR} \\
\texttt{SOLVE: ```python ... ```} \\
            }
        }
    }
     \\ \\

    \theutterance \stepcounter{utterance}  
    & \multicolumn{4}{p{0.6\linewidth}}{
        \cellcolor[rgb]{0.561,0.737,0.561}{
            \makecell[{{p{\linewidth}}}]{
                \texttt{\tiny{[P1$\rangle$GM]}}
                \texttt{\{"current\_mode":"don't know","input\_rationale":"The only example preserves a string of spaces exactly, which is most consistent with the identity function. With no other examples, this is my best commit.","hypothesis":"def f(text: str) {-}> str:\textbackslash{}n    \textbackslash{}"\textbackslash{}"\textbackslash{}"Returns the input text unchanged.\textbackslash{}"\textbackslash{}"\textbackslash{}"\textbackslash{}n    return text","confidence\_score":62\}} \\
            }
        }
    }
    & & \\ \\

    \theutterance \stepcounter{utterance}  
    & & & \multicolumn{4}{p{0.6\linewidth}}{
        \cellcolor[rgb]{0.9,0.9,0.9}{
            \makecell[{{p{\linewidth}}}]{
                \texttt{\tiny{[P1$\langle$GM]}}
                \texttt{OUTPUT: ('a{-}b', 'n{-}o')} \\
            }
        }
    }
     \\ \\

    \theutterance \stepcounter{utterance}  
    & \multicolumn{4}{p{0.6\linewidth}}{
        \cellcolor[rgb]{0.561,0.737,0.561}{
            \makecell[{{p{\linewidth}}}]{
                \texttt{\tiny{[P1$\rangle$GM]}}
                \texttt{\{"current\_mode":"don't know","input\_rationale":"The example 'a{-}b' {-}> 'n{-}o' matches a ROT13{-}style Caesar shift by 13 on letters, preserving punctuation and spaces. The earlier spaces{-}only example is also consistent with this.","hypothesis":"def f(text: str) {-}> str:\textbackslash{}n    \textbackslash{}"\textbackslash{}"\textbackslash{}"Applies ROT13 to letters, leaving all other characters unchanged.\textbackslash{}"\textbackslash{}"\textbackslash{}"\textbackslash{}n    out = {[}{]}\textbackslash{}n    for ch in text:\textbackslash{}n        o = ord(ch)\textbackslash{}n        if 65 <= o <= 90:\textbackslash{}n            out.append(chr((o {-} 65 + 13) \% 26 + 65))\textbackslash{}n        elif 97 <= o <= 122:\textbackslash{}n            out.append(chr((o {-} 97 + 13) \% 26 + 97))\textbackslash{}n        else:\textbackslash{}n            out.append(ch)\textbackslash{}n    return ''.join(out)","confidence\_score":99\}} \\
            }
        }
    }
    & & \\ \\

    \theutterance \stepcounter{utterance}  
    & & & \multicolumn{4}{p{0.6\linewidth}}{
        \cellcolor[rgb]{0.9,0.9,0.9}{
            \makecell[{{p{\linewidth}}}]{
                \texttt{\tiny{[P1$\langle$GM]}}
                \texttt{OUTPUT: ('a', 'n')} \\
            }
        }
    }
     \\ \\

    \theutterance \stepcounter{utterance}  
    & \multicolumn{4}{p{0.6\linewidth}}{
        \cellcolor[rgb]{0.561,0.737,0.561}{
            \makecell[{{p{\linewidth}}}]{
                \texttt{\tiny{[P1$\rangle$GM]}}
                \texttt{\{"current\_mode":"don't know","input\_rationale":"The new example 'a' {-}> 'n' confirms a ROT13 shift on letters, with non{-}letters preserved.","hypothesis":"def f(text: str) {-}> str:\textbackslash{}n    \textbackslash{}"\textbackslash{}"\textbackslash{}"Applies ROT13 to letters, leaving all other characters unchanged.\textbackslash{}"\textbackslash{}"\textbackslash{}"\textbackslash{}n    out = {[}{]}\textbackslash{}n    for ch in text:\textbackslash{}n        o = ord(ch)\textbackslash{}n        if 65 <= o <= 90:\textbackslash{}n            out.append(chr((o {-} 65 + 13) \% 26 + 65))\textbackslash{}n        elif 97 <= o <= 122:\textbackslash{}n            out.append(chr((o {-} 97 + 13) \% 26 + 97))\textbackslash{}n        else:\textbackslash{}n            out.append(ch)\textbackslash{}n    return ''.join(out)","confidence\_score":100\}} \\
            }
        }
    }
    & & \\ \\

    \theutterance \stepcounter{utterance}  
    & & & \multicolumn{4}{p{0.6\linewidth}}{
        \cellcolor[rgb]{0.9,0.9,0.9}{
            \makecell[{{p{\linewidth}}}]{
                \texttt{\tiny{[P1$\langle$GM]}}
                \texttt{OUTPUT: ('aba', 'non')} \\
            }
        }
    }
     \\ \\

    \theutterance \stepcounter{utterance}  
    & \multicolumn{4}{p{0.6\linewidth}}{
        \cellcolor[rgb]{0.561,0.737,0.561}{
            \makecell[{{p{\linewidth}}}]{
                \texttt{\tiny{[P1$\rangle$GM]}}
                \texttt{\{"current\_mode":"don't know","input\_rationale":"The example 'aba' {-}> 'non' is exactly ROT13 applied character by character, confirming the hypothesis with repeated letters as well.","hypothesis":"def f(text: str) {-}> str:\textbackslash{}n    \textbackslash{}"\textbackslash{}"\textbackslash{}"Applies ROT13 to letters, leaving all other characters unchanged.\textbackslash{}"\textbackslash{}"\textbackslash{}"\textbackslash{}n    out = {[}{]}\textbackslash{}n    for ch in text:\textbackslash{}n        o = ord(ch)\textbackslash{}n        if 65 <= o <= 90:\textbackslash{}n            out.append(chr((o {-} 65 + 13) \% 26 + 65))\textbackslash{}n        elif 97 <= o <= 122:\textbackslash{}n            out.append(chr((o {-} 97 + 13) \% 26 + 97))\textbackslash{}n        else:\textbackslash{}n            out.append(ch)\textbackslash{}n    return ''.join(out)","confidence\_score":100\}} \\
            }
        }
    }
    & & \\ \\

    \theutterance \stepcounter{utterance}  
    & & & \multicolumn{4}{p{0.6\linewidth}}{
        \cellcolor[rgb]{0.9,0.9,0.9}{
            \makecell[{{p{\linewidth}}}]{
                \texttt{\tiny{[P1$\langle$GM]}}
                \texttt{OUTPUT: (' ', ' ')} \\
            }
        }
    }
     \\ \\

    \theutterance \stepcounter{utterance}  
    & \multicolumn{4}{p{0.6\linewidth}}{
        \cellcolor[rgb]{0.561,0.737,0.561}{
            \makecell[{{p{\linewidth}}}]{
                \texttt{\tiny{[P1$\rangle$GM]}}
                \texttt{\{"current\_mode":"don't know","input\_rationale":"The example with a single space is also consistent with ROT13, since non{-}letters are preserved unchanged.","hypothesis":"def f(text: str) {-}> str:\textbackslash{}n    \textbackslash{}"\textbackslash{}"\textbackslash{}"Applies ROT13 to letters, leaving all other characters unchanged.\textbackslash{}"\textbackslash{}"\textbackslash{}"\textbackslash{}n    out = {[}{]}\textbackslash{}n    for ch in text:\textbackslash{}n        o = ord(ch)\textbackslash{}n        if 65 <= o <= 90:\textbackslash{}n            out.append(chr((o {-} 65 + 13) \% 26 + 65))\textbackslash{}n        elif 97 <= o <= 122:\textbackslash{}n            out.append(chr((o {-} 97 + 13) \% 26 + 97))\textbackslash{}n        else:\textbackslash{}n            out.append(ch)\textbackslash{}n    return ''.join(out)","confidence\_score":100\}} \\
            }
        }
    }
    & & \\ \\

    \theutterance \stepcounter{utterance}  
    & & & \multicolumn{4}{p{0.6\linewidth}}{
        \cellcolor[rgb]{0.9,0.9,0.9}{
            \makecell[{{p{\linewidth}}}]{
                \texttt{\tiny{[P1$\langle$GM]}}
                \texttt{OUTPUT: ('A', 'A')} \\
            }
        }
    }
     \\ \\

    \theutterance \stepcounter{utterance}  
    & \multicolumn{4}{p{0.6\linewidth}}{
        \cellcolor[rgb]{0.561,0.737,0.561}{
            \makecell[{{p{\linewidth}}}]{
                \texttt{\tiny{[P1$\rangle$GM]}}
                \texttt{\{"current\_mode":"don't know","input\_rationale":"This confirms the function leaves uppercase A unchanged, so the transform is not ROT13. Combined with prior examples, it appears to return strings unchanged.","hypothesis":"def f(text: str) {-}> str:\textbackslash{}n    \textbackslash{}"\textbackslash{}"\textbackslash{}"Returns the input text unchanged.\textbackslash{}"\textbackslash{}"\textbackslash{}"\textbackslash{}n    return text","confidence\_score":99\}} \\
            }
        }
    }
    & & \\ \\

    \theutterance \stepcounter{utterance}  
    & & & \multicolumn{4}{p{0.6\linewidth}}{
        \cellcolor[rgb]{0.9,0.9,0.9}{
            \makecell[{{p{\linewidth}}}]{
                \texttt{\tiny{[P1$\langle$GM]}}
                \texttt{OUTPUT: ('123', '123')} \\
            }
        }
    }
     \\ \\

    \theutterance \stepcounter{utterance}  
    & \multicolumn{4}{p{0.6\linewidth}}{
        \cellcolor[rgb]{0.561,0.737,0.561}{
            \makecell[{{p{\linewidth}}}]{
                \texttt{\tiny{[P1$\rangle$GM]}}
                \texttt{\{"current\_mode":"don't know","input\_rationale":"The examples now show spaces, letters, and digits all preserved exactly, which strongly indicates the function is the identity transform on the input string.","hypothesis":"def f(text: str) {-}> str:\textbackslash{}n    \textbackslash{}"\textbackslash{}"\textbackslash{}"Returns the input text unchanged.\textbackslash{}"\textbackslash{}"\textbackslash{}"\textbackslash{}n    return text","confidence\_score":100\}} \\
            }
        }
    }
    & & \\ \\

    \theutterance \stepcounter{utterance}  
    & & & \multicolumn{4}{p{0.6\linewidth}}{
        \cellcolor[rgb]{0.9,0.9,0.9}{
            \makecell[{{p{\linewidth}}}]{
                \texttt{\tiny{[P1$\langle$GM]}}
                \texttt{OUTPUT: ('', '')} \\
            }
        }
    }
     \\ \\

    \theutterance \stepcounter{utterance}  
    & \multicolumn{4}{p{0.6\linewidth}}{
        \cellcolor[rgb]{0.561,0.737,0.561}{
            \makecell[{{p{\linewidth}}}]{
                \texttt{\tiny{[P1$\rangle$GM]}}
                \texttt{\{"current\_mode":"don't know","input\_rationale":"The empty{-}string example is also consistent with the identity transform, which is now confirmed by all examples across letters, digits, spaces, and punctuation{-}like cases.","hypothesis":"def f(text: str) {-}> str:\textbackslash{}n    \textbackslash{}"\textbackslash{}"\textbackslash{}"Returns the input text unchanged.\textbackslash{}"\textbackslash{}"\textbackslash{}"\textbackslash{}n    return text","confidence\_score":100\}} \\
            }
        }
    }
    & & \\ \\

    \theutterance \stepcounter{utterance}  
    & & & \multicolumn{4}{p{0.6\linewidth}}{
        \cellcolor[rgb]{0.9,0.9,0.9}{
            \makecell[{{p{\linewidth}}}]{
                \texttt{\tiny{[P1$\langle$GM]}}
                \texttt{OUTPUT: ('aA', 'nA')} \\
            }
        }
    }
     \\ \\

    \theutterance \stepcounter{utterance}  
    & \multicolumn{4}{p{0.6\linewidth}}{
        \cellcolor[rgb]{0.561,0.737,0.561}{
            \makecell[{{p{\linewidth}}}]{
                \texttt{\tiny{[P1$\rangle$GM]}}
                \texttt{\{"current\_mode":"don't know","input\_rationale":"The new example 'aA' {-}> 'nA' shows lowercase letters are transformed by ROT13 while uppercase letters are left unchanged. Combined with earlier examples, the function applies ROT13 only to lowercase ASCII letters and preserves everything else.","hypothesis":"def f(text: str) {-}> str:\textbackslash{}n    \textbackslash{}"\textbackslash{}"\textbackslash{}"Applies ROT13 to lowercase letters only, leaving all other characters unchanged.\textbackslash{}"\textbackslash{}"\textbackslash{}"\textbackslash{}n    out = {[}{]}\textbackslash{}n    for ch in text:\textbackslash{}n        o = ord(ch)\textbackslash{}n        if 97 <= o <= 122:\textbackslash{}n            out.append(chr((o {-} 97 + 13) \% 26 + 97))\textbackslash{}n        else:\textbackslash{}n            out.append(ch)\textbackslash{}n    return ''.join(out)","confidence\_score":100\}} \\
            }
        }
    }
    & & \\ \\

    \theutterance \stepcounter{utterance}  
    & & & \multicolumn{4}{p{0.6\linewidth}}{
        \cellcolor[rgb]{0.9,0.9,0.9}{
            \makecell[{{p{\linewidth}}}]{
                \texttt{\tiny{[P1$\langle$GM]}}
                \texttt{OUTPUT: ('hello', 'uryyb')} \\
            }
        }
    }
     \\ \\

    \theutterance \stepcounter{utterance}  
    & \multicolumn{4}{p{0.6\linewidth}}{
        \cellcolor[rgb]{0.561,0.737,0.561}{
            \makecell[{{p{\linewidth}}}]{
                \texttt{\tiny{[P1$\rangle$GM]}}
                \texttt{\{"current\_mode":"don't know","input\_rationale":"The latest example 'hello' {-}> 'uryyb' confirms lowercase ROT13. Earlier examples show uppercase letters are left unchanged and non{-}letters are preserved, so the function applies ROT13 only to lowercase ASCII letters.","hypothesis":"def f(text: str) {-}> str:\textbackslash{}n    \textbackslash{}"\textbackslash{}"\textbackslash{}"Applies ROT13 to lowercase letters only, leaving all other characters unchanged.\textbackslash{}"\textbackslash{}"\textbackslash{}"\textbackslash{}n    out = {[}{]}\textbackslash{}n    for ch in text:\textbackslash{}n        o = ord(ch)\textbackslash{}n        if 97 <= o <= 122:\textbackslash{}n            out.append(chr((o {-} 97 + 13) \% 26 + 97))\textbackslash{}n        else:\textbackslash{}n            out.append(ch)\textbackslash{}n    return ''.join(out)","confidence\_score":100\}} \\
            }
        }
    }
    & & \\ \\

    \theutterance \stepcounter{utterance}  
    & & & \multicolumn{4}{p{0.6\linewidth}}{
        \cellcolor[rgb]{0.9,0.9,0.9}{
            \makecell[{{p{\linewidth}}}]{
                \texttt{\tiny{[P1$\langle$GM]}}
                \texttt{OUTPUT: ('ab', 'no')} \\
            }
        }
    }
     \\ \\

    \theutterance \stepcounter{utterance}  
    & \multicolumn{4}{p{0.6\linewidth}}{
        \cellcolor[rgb]{0.561,0.737,0.561}{
            \makecell[{{p{\linewidth}}}]{
                \texttt{\tiny{[P1$\rangle$GM]}}
                \texttt{\{"current\_mode":"don't know","input\_rationale":"The example 'ab' {-}> 'no' matches lowercase ROT13 exactly, consistent with all prior observations that only lowercase ASCII letters are shifted and everything else is preserved.","hypothesis":"def f(text: str) {-}> str:\textbackslash{}n    \textbackslash{}"\textbackslash{}"\textbackslash{}"Applies ROT13 to lowercase letters only, leaving all other characters unchanged.\textbackslash{}"\textbackslash{}"\textbackslash{}"\textbackslash{}n    out = {[}{]}\textbackslash{}n    for ch in text:\textbackslash{}n        o = ord(ch)\textbackslash{}n        if 97 <= o <= 122:\textbackslash{}n            out.append(chr((o {-} 97 + 13) \% 26 + 97))\textbackslash{}n        else:\textbackslash{}n            out.append(ch)\textbackslash{}n    return ''.join(out)","confidence\_score":100\}} \\
            }
        }
    }
    & & \\ \\

    \theutterance \stepcounter{utterance}  
    & & & \multicolumn{4}{p{0.6\linewidth}}{
        \cellcolor[rgb]{0.9,0.9,0.9}{
            \makecell[{{p{\linewidth}}}]{
                \texttt{\tiny{[P1$\langle$GM]}}
                \texttt{OUTPUT: ('hello world', 'uryyb jbeyq')} \\
            }
        }
    }
     \\ \\

    \theutterance \stepcounter{utterance}  
    & \multicolumn{4}{p{0.6\linewidth}}{
        \cellcolor[rgb]{0.561,0.737,0.561}{
            \makecell[{{p{\linewidth}}}]{
                \texttt{\tiny{[P1$\rangle$GM]}}
                \texttt{\{"current\_mode":"don't know","input\_rationale":"The example 'hello world' {-}> 'uryyb jbeyq' confirms lowercase ROT13 across words and preserves spaces, matching the current hypothesis exactly.","hypothesis":"def f(text: str) {-}> str:\textbackslash{}n    \textbackslash{}"\textbackslash{}"\textbackslash{}"Applies ROT13 to lowercase letters only, leaving all other characters unchanged.\textbackslash{}"\textbackslash{}"\textbackslash{}"\textbackslash{}n    out = {[}{]}\textbackslash{}n    for ch in text:\textbackslash{}n        o = ord(ch)\textbackslash{}n        if 97 <= o <= 122:\textbackslash{}n            out.append(chr((o {-} 97 + 13) \% 26 + 97))\textbackslash{}n        else:\textbackslash{}n            out.append(ch)\textbackslash{}n    return ''.join(out)","confidence\_score":100\}} \\
            }
        }
    }
    & & \\ \\

    \theutterance \stepcounter{utterance}  
    & & & \multicolumn{4}{p{0.6\linewidth}}{
        \cellcolor[rgb]{0.9,0.9,0.9}{
            \makecell[{{p{\linewidth}}}]{
                \texttt{\tiny{[P1$\langle$GM]}}
                \texttt{OUTPUT: ('\_kM', '\_xM')} \\
            }
        }
    }
     \\ \\

    \theutterance \stepcounter{utterance}  
    & \multicolumn{4}{p{0.6\linewidth}}{
        \cellcolor[rgb]{0.561,0.737,0.561}{
            \makecell[{{p{\linewidth}}}]{
                \texttt{\tiny{[P1$\rangle$GM]}}
                \texttt{\{"current\_mode":"don't know","input\_rationale":"The example '\_kM' {-}> '\_xM' confirms the same rule: lowercase letters are ROT13{-}shifted, while uppercase letters and non{-}letters are left unchanged.","hypothesis":"def f(text: str) {-}> str:\textbackslash{}n    \textbackslash{}"\textbackslash{}"\textbackslash{}"Applies ROT13 to lowercase letters only, leaving all other characters unchanged.\textbackslash{}"\textbackslash{}"\textbackslash{}"\textbackslash{}n    out = {[}{]}\textbackslash{}n    for ch in text:\textbackslash{}n        o = ord(ch)\textbackslash{}n        if 97 <= o <= 122:\textbackslash{}n            out.append(chr((o {-} 97 + 13) \% 26 + 97))\textbackslash{}n        else:\textbackslash{}n            out.append(ch)\textbackslash{}n    return ''.join(out)","confidence\_score":100\}} \\
            }
        }
    }
    & & \\ \\

    \theutterance \stepcounter{utterance}  
    & & & \multicolumn{4}{p{0.6\linewidth}}{
        \cellcolor[rgb]{0.9,0.9,0.9}{
            \makecell[{{p{\linewidth}}}]{
                \texttt{\tiny{[P1$\langle$GM]}}
                \texttt{OUTPUT: ('ts', 'gf')} \\
            }
        }
    }
     \\ \\

    \theutterance \stepcounter{utterance}  
    & \multicolumn{4}{p{0.6\linewidth}}{
        \cellcolor[rgb]{0.561,0.737,0.561}{
            \makecell[{{p{\linewidth}}}]{
                \texttt{\tiny{[P1$\rangle$GM]}}
                \texttt{\{"current\_mode":"don't know","input\_rationale":"The example 'ts' {-}> 'gf' is consistent with lowercase ROT13, reinforcing the same rule across arbitrary lowercase input.","hypothesis":"def f(text: str) {-}> str:\textbackslash{}n    \textbackslash{}"\textbackslash{}"\textbackslash{}"Applies ROT13 to lowercase letters only, leaving all other characters unchanged.\textbackslash{}"\textbackslash{}"\textbackslash{}"\textbackslash{}n    out = {[}{]}\textbackslash{}n    for ch in text:\textbackslash{}n        o = ord(ch)\textbackslash{}n        if 97 <= o <= 122:\textbackslash{}n            out.append(chr((o {-} 97 + 13) \% 26 + 97))\textbackslash{}n        else:\textbackslash{}n            out.append(ch)\textbackslash{}n    return ''.join(out)","confidence\_score":100\}} \\
            }
        }
    }
    & & \\ \\

    \theutterance \stepcounter{utterance}  
    & & & \multicolumn{2}{p{0.3\linewidth}}{
        \cellcolor[rgb]{0.9,0.9,0.9}{
            \makecell[{{p{\linewidth}}}]{
                \texttt{\tiny{[GM$|$GM]}}
                \texttt{Game Verdict: LOSS (turn limit reached)} \\
            }
        }
    }
    & & \\ \\

\end{supertabular}
}

\end{document}
